\section{Related Work}
\label{sec:related}

\subsection{LLM Benchmarks for EDA and Hardware Design}

\textbf{RTL code generation.} RTLLM~\cite{lu2024rtllm} introduced 30 (later 50) digital designs with natural-language specifications, testbenches, and reference RTL, evaluating syntax correctness, functional accuracy, and PPA quality. VerilogEval~\cite{liu2023verilogeval} extended this to 100+ Verilog tasks from diverse sources. Both evaluate \emph{single-module} generation---one function, one testbench---without tool interaction or multi-step orchestration. \bench{} differs fundamentally: each task requires the agent to call live EDA tools, inspect outputs, and produce a multi-file deliverable across multiple steps.

\textbf{Flow repair.} FlowFixer-EDA~\cite{dong2026flowfixer} and EDA-FlowBench evaluate diagnosing and repairing failures in open-source RTL-to-GDS flows (Yosys, OpenROAD), with 200 tasks covering multiple error categories. While closer to real workflows, these benchmarks operate on \emph{open-source tools only} and evaluate \emph{repair} (fixing a broken run) rather than \emph{orchestration} (building a deliverable from scratch). \bench{} uses commercial EDA tools (Empyrean Aether, Cadence Virtuoso, Synopsys tools) and evaluates end-to-end deliverable production.

\textbf{EDA scripting.} EDA-Scripting-Bench~\cite{zhu2026edascript} evaluates LLM agents on PyAether, SKILL, and Tcl scripting across 449 tasks with execution-based scoring. While sharing the execution-grounded philosophy, it evaluates \emph{single-script} generation (one function per task), not multi-step flow orchestration with deliverable packaging. \bench{} extends this to the \mode{} paradigm where the agent must produce a complete deliverable package.

\subsection{Agent Benchmarks in Software Engineering}

SWE-bench~\cite{jimenez2024swebench} and its variants (SWE-bench Verified, SWE-bench Pro) evaluate whether AI agents can resolve real GitHub issues by editing codebases. The key insight---execution-based evaluation against a project's own test suite---inspires our design. However, SWE-bench operates on Python repositories with Docker-based evaluation, while \bench{} operates on EDA design flows with live commercial tool execution on a remote EDA server. The EDA setting introduces unique challenges: closed tools, non-redistributable data, and tool-chain grounding requirements that have no software-engineering analog.

\subsection{Tool-Use and Execution-Grounded Agents}

The MCP (Model Context Protocol)~\cite{lee2025autoeda} standardizes tool exposure to LLMs. Execution-grounded agents~\cite{zhu2026edascript, schick2024toolformer} close a generate--execute--refine loop by observing tool outputs. \bench{} evaluates agents in this paradigm: the agent must call EDA tools through an MCP bridge, read execution results (simulation logs, DRC reports, netlist extracts), and iteratively refine its deliverable. This is qualitatively harder than single-shot generation because the agent must maintain state across tool calls, interpret domain-specific outputs, and decide when the deliverable is complete.
